\subsection{Možnost kombinace přístupů}
\label{subsec:static-llm-combination}

Z výše uvedeného srovnání vyplývá, že statické analyzátory a LLM by neměly být vnímány jako konkurenční, ale jako komplementární technologie.
Každý z přístupů pokrývá oblast, kde druhý selhává nebo je slabší.

Statická analýza je silná tam, kde je třeba formální, reprodukovatelná a úplná detekce přesně definovaných vzorů, jako jsou syntaktické chyby, jednoduché úniky paměti či porušení kódovacích standardů.
LLM je naopak silný tam, kde je třeba kontextová interpretace, srozumitelné vysvětlení a hodnocení aspektů, které nejsou formalizovatelné jako pravidla.
Mezi ně patří posouzení adekvátnosti algoritmu, celkového designu řešení, koherence pojmenování nebo pochopení záměru zadání.

V prostředí systému Kelvin je vhodné uvažovat o hybridním přístupu, kde automatické testy a statická analýza pokrývají vrstvu formální správnosti, zatímco LLM přidává vrstvu kontextového hodnocení a srozumitelné zpětné vazby.
Výstupy statické analýzy mohou být navíc zahrnuty jako součást vstupu do LLM: model může být požádán, aby vysvětlil konkrétní upozornění \texttt{cppcheck} nebo navrhl způsob jeho opravy v kontextu celého řešení.
Tímto způsobem LLM nezdvojuje práci statického analyzátoru, ale přidává vrstvu interpretace, která formálnímu nástroji chybí.

Je však nutné zachovat realistická očekávání, kde ani kombinace obou přístupů nenahrazuje odborný úsudek vyučujícího.
Statická analýza ani LLM nemohou zaručit, že odevzdané řešení je správné, efektivní nebo didakticky přínosné.
Slouží jako podpůrné nástroje usnadňující orientaci v řešení, \textbf{nikoli jako jeho hodnotitel}.

\endinput